% =====================================================================
%  tr03b —— 原文 pp.23–25
%  内容：§3.3 奇异轨迹（续）、§3.4 问题集 II（问题 6–10）
%  编号：问题 6、7、8、9、10（纯数字，\begin{prob}{N}）
%  边界：p.23 开头承接 tr03a 末段（跨页句已在 tr03a 写完），此处从新段
%        "The most singular fiber …" 起；终点 p.25 末句 "…Hitchin maps defined
%        in [72] and [10]."（p.26 由 tr04 从 §4 接）。
% =====================================================================

% [原文 p.23]

最奇异的纤维是 $0\in\cA_{G_{\C}}$ 上方的纤维，Laumon [71] 称之为幂零锥（nilpotent cone），以强调它与李代数中幂零锥的类比。这一纤维已从不同角度得到研究，其中之一是通过 $S^1$ 作用 $(E,\Phi)\to(E,\ee^{\ii\theta}\Phi)$ 的矩映射（moment map）$\mu$ 来研究，如 [57,60] 所做。特别地，注意：给定稳定丛 $E$，取 $\Phi=0$，则 Higgs 丛 $(E,0)$ 是稳定的，并且被 $S^1$-作用平凡地固定。

% [Tutor Note: p.23 把 $0\in\cA_{G_{\C}}$ 上方的纤维称为幂零锥，并称之为"最奇异"的纤维；由 Hitchin 映射 $h$ 的定义，该纤维即 $h^{-1}(0)$，它包含所有 $(E,0)$（$E$ 稳定），这正是紧接的"$(E,0)$ 稳定且被 $S^1$ 作用平凡固定"的含义，也是 $S^1$ 矩映射流把 $\cM_{G_{\C}}$ 中诸点吸引过去的原因。另注意原文此处映射记号写作 $(E,\Phi)\to(E,\ee^{\ii\theta}\Phi)$，实际含义是按 $\Phi\mapsto\ee^{\ii\theta}\Phi$ 作用。]

如 [56, Theorem 5.2] 所示，幂零锥在由 $\mu$ 生成的流下保持不变，并且它编码了模空间的拓扑，因为 $\cM_{G_{\C}}$ 中的点会沿流趋向幂零锥。关于奇异 Higgs 丛以及 Hitchin 纤维化奇异轨迹上方的 Higgs 丛的研究概况，读者可参阅 [6]。特别地，幂零锥主要是就 $\SL(n,\C)$ 与 $\GL(n,\C)$ 进行研究的，而对于 $G_{\C}$-Higgs 丛的模空间，其几何的很大一部分仍属未知。对 $\SL(n,\C)$ 与 $\GL(n,\C)$-Higgs 丛，幂零锥的不可约分支由 $S^1$ 作用的不动点集的连通分支所标记。幂零锥也可以对其他类型的 Higgs 丛来考虑，我们将在第 5 节中就通过超多边形（hyperpolygons）出现的幂零锥略作评论。在经典 Higgs 丛的幂零锥的各分支中，包含半稳定丛的模空间 $\cN$。

Hitchin 纤维化奇异轨迹上方的其他纤维则是近年研究主题（例如见 [11,30,52,64,92,97]）。由于我们将在研究 Higgs 丛的膜（brane）时回到奇异纤维，这里以几点简短评论结束本章：

\begin{itemize}[leftmargin=1.4em,itemsep=1pt,topsep=1pt]
\item 对 $\GL(n,\C)$-Higgs 丛，我们提到过：对 Hitchin 底的一般点，相应的纤维可以等同于谱曲线 $S$ 的雅可比簇 $\Jac(S)$。当谱曲线 $S$ 不光滑但整的（integral）时，相应的纤维可视为紧化雅可比簇（compactified Jacobian）[18,94--96]（关于其清楚的解释，另见 [75, Fact 10.3]）。紧化雅可比簇 $\overline{\Jac}(S)$ 是 $S$ 上所有无挠秩 1 层的模空间，其中通常的“局部自由”条件被去掉了。直观地说，紧化雅可比簇可以通过考虑一族趋近于奇异曲线 $S_0$ 的光滑曲线 $S_t$ 来得到，其中 $S_0$ 是幂零锥的基点；由于 $\Jac(S_t)$ 的极限不依赖于光滑族的选取 [68]，该极限即为 $\overline{\Jac}(S)$。
\item 对谱曲线 $S$ 不整的 $\GL(n,\C)$-Higgs 丛，需要考虑精细模空间（fine moduli space）。
\item 对 $\SL(2,\C)$-Higgs 丛，当 $S$ 不可约且仅有简单结点时，$\cM_{\Sl(2,\C)}$ 的纤维的连通性已在 [45, §5.2.2] 中考虑，而奇异纤维的完整描述见 [52]。
\end{itemize}

% [Tutor Note: p.23 的 "integral"（整的）指曲线既约且不可约；当谱曲线 $S$ 整而奇异时，Hitchin 纤维同构于紧化雅可比簇 $\overline{\Jac}(S)$，即 $S$ 上无挠秩 1 层的模空间（这是模空间理论中的经典结果，见原文所引 [18,94--96]）；只有在 $S$ 光滑时它才等同于 $\Jac(S)$。]

% [原文 p.24]

\subsection{问题集 II}

本问题集旨在让读者熟悉复 Higgs 丛的谱数据，并理解它如何能以阿贝尔的方式表达。此外，它应当为研究 Hitchin 纤维化奇异纤维的种种方式提供一个入门。这些问题既有来自讲义的问题，也有一些补充问题。

\begin{prob}{6}
\textbf{谱曲线。}通过其特征多项式，Higgs 场在典范丛 $K$ 的全空间中定义了一条自然的谱曲线。此外，Higgs 场所满足的额外条件尤其体现在这条曲线的性质之中。

\begin{enumerate}[label=(\roman*),leftmargin=2.2em]
\item 在 $A$、$B$、$C$、$D$ 型复李群中，哪些群所对应的 Higgs 场所定义的谱曲线带有对合？
\item 考虑低秩群，找出 Higgs 丛 $(E,\Phi)$ 所应满足的条件，使得该 Higgs 场定义的谱曲线具有一个 4 阶自同构（例如见 [13]）。
\item 从线性代数的角度考察特征多项式的性质，给出两类特殊 Higgs 丛 $(E,\Phi)$ 的描述，使得谱曲线 $S=\{(\det\Phi-\lambda)=0\}$ 的方程为 $\{q^3(\lambda)=0\}$，其中 $q(\lambda)$ 为某个不可约多项式。
\item 承问题 6.iii.，对哪些 Higgs 丛有
\[
(\det\Phi-\lambda)=q^n(\lambda)
\]
其中 $q(\lambda)$ 为某个不可约多项式？
\end{enumerate}
\end{prob}

% [Tutor Note: p.24 问题 6.iii 中谱曲线写作集合 $S=\{(\det\Phi-\lambda)=0\}$、其方程写作 $\{q^3(\lambda)=0\}$，这里的外层花括号是原文的集合记号（不属于方程本身）；实际含义是"$\det\Phi-\lambda=0$ 的零点轨迹恰由 $q^3(\lambda)=0$ 给出"，即特征多项式是某个不可约多项式的三次幂。]

\begin{prob}{7}
\textbf{一般纤维。}我们已经提到，通过对 Higgs 对 $(E,\Phi)$ 添加条件，可以得到主 $G_{\C}$-Higgs 丛的描述。等价地，也可以对经典谱曲线 $S$ 及其上的数据 $\Jac(S)$ 添加条件，从而还原主 Higgs 丛。

\begin{enumerate}[label=(\roman*),leftmargin=2.2em]
\item 在考虑经典 Higgs 丛的谱数据时，利用 Grothendieck-Riemann-Roch 证明 $\deg(E)=\deg(L)+(n^2-n)(1-g)$。
\item 证明：在考虑 $\SL(n,\C)$-Higgs 丛时，条件
\[
\Lambda^nE\cong\cO
\]
等价于要求谱线丛 $L$ 属于 $\Prym(S,\Sigma)$。
\item 证明：光滑的 $\SO(2n,\C)$-Hitchin 纤维是 $\Prym(\hat S,\hat S/\hat\sigma)$，其中 $\hat S$ 是与正则基点相对应的曲线 $S$ 的消解（desingularization）。
\item 证明：在考虑一般的 $\SO(2n+1,\C)$-Higgs 丛时，由特征多项式定义的谱曲线 $S$ 绝不光滑。
\end{enumerate}
\end{prob}

% [Tutor Note: p.24 问题 7.i 的公式可用 $\chi(L)=\chi(\pi_*L)$ 验算：在 $S$ 上 $\chi(L)=\deg L+1-g_S$，而 $\pi_*L$ 是 $\Sigma$ 上秩 $n$ 的丛，$\chi(\pi_*L)=\deg(\pi_*L)+n(1-g)$；代入 $g_S=n^2(g-1)+1$ 即得 $\deg E=\deg L+(n^2-n)(1-g)$。原句 "Use Grothendieck-Riemann-Roch to $\deg(E)=\cdots$" 脱了动词（应为 to show/to derive）。]

% [Tutor Note: p.24 问题 7.iii 中 $\Prym(\hat S,\hat S/\hat\sigma)$ 是 $2:1$ 覆盖 $\hat S\to\hat S/\hat\sigma$ 的 Prym 簇：第二个参数取商曲线（而不是 $\Sigma$），说明 $\hat\sigma$ 是 $\hat S$ 上由 $\sigma$ 诱导的对合、$\hat S$ 是与之相应的消解曲线。]

\begin{prob}{8}
\textbf{Hitchin 纤维化。}当考虑与正则轨迹只有有限交的 Higgs 丛子空间时，人们可以通过单值性作用以及其他对合来研究其几何与拓扑。

% [Tutor Note: p.24 问题 8 正文 "…study its geometry and topology though the monodromy action…" 中 "though" 应为 "through"（原文笔误）。]

% [原文 p.25]
\begin{enumerate}[label=(\roman*),leftmargin=2.2em]
\item 李代数 $\fg_{\C}$ 上有一个自然对合 $\sigma_s$，由
\[
\sigma_s([\operatorname{ad}f_0]^n e_i)=(-1)^{n+1}[\operatorname{ad}f_0]^n e_i
\]
给出。该自同构由 $\sigma_s(e_i)=-e_i$ 与 $\sigma_s(f_0)=-f_0$ 唯一确定。证明：$\fg_{\C}$ 上的 $-\sigma_s$ 在 $\fg_{\C}$ 的不变多项式环上平凡地作用。
\item 证明：对非分裂实形式 $G$，$G$-Higgs 场的特征多项式定义一条可约曲线。
\item ((*)) 推广 [19]，在 $X$ 上定义广义抛物 $G_{\C}$-Higgs 丛以及相应的 Hitchin 纤维化。
\end{enumerate}
\end{prob}

% [Tutor Note: p.25 问题 8.i 的 $\sigma_s$ 即分裂（split）实形式所对应的对合：由 $\sigma_s(e_i)=-e_i$、$\sigma_s(f_0)=-f_0$ 及 $\sigma_s$ 是李代数自同构，可归纳验证 $\sigma_s([\operatorname{ad}f_0]^n e_i)=[\operatorname{ad}(-f_0)]^n\sigma_s(e_i)=(-1)^{n+1}[\operatorname{ad}f_0]^ne_i$（$n=0$ 时即 $\sigma_s(e_i)=-e_i$，故 $n+1$ 次幂的符号正确）。]

\begin{prob}{9}
\textbf{奇异纤维。}考虑 $\SU(p,p+q)$-Higgs 丛 $(V\oplus W,\Phi)$，其特征多项式为
\[
\det(\eta-\Phi)=\eta^q(\eta^{2p}+a_2\eta^{2p-2}+\cdots+a_{2p-2}\eta^2+a_{2p}),
\]
并带有自然对合 $\sigma:\eta\mapsto-\eta$，如 (3.9) 所述。

\begin{enumerate}[label=(\roman*),leftmargin=2.2em]
\item 当 $q=0$ 时，考虑
\[
\bar\pi:\bar S=S/\sigma\to\Sigma
\]
上的线丛 $U_1$ 与 $U_2$，使得 $\bar\pi_*U_1=V$ 且 $\bar\pi_*U_1=W$，并用 $U_i$ 表示 $\Lambda^pV$、$\Lambda^pW$ 以及 $\Lambda^pV\cong\Lambda^pW^*$。
\item 当 $p=2$、$q=0$ 时，问题 9.i. 中的商曲线 $\bar S$ 是黎曼面 $\Sigma$ 的二重覆盖。描述由该 Higgs 丛的秩较低这一事实所出现的任何额外几何结构。
\item ((*)) $\SO(2,4)$-Higgs 丛的结构如何能通过问题 9.ii. 所描述的数据呈现出来？
\item ((*)) 考虑 $p=2$、$q=2$，给出可与 $\U(2,4)$-Higgs 丛相联系的一种谱数据描述--关于带有符号差的正交 Higgs 丛的谱数据，细节见 [11]。
\end{enumerate}
\end{prob}

% [Tutor Note: p.25 问题 9.i 原文作 "$\bar\pi_*U_1=V$ and $\bar\pi_*U_1=W$"，两个直像都用了 $U_1$；由题设（$U_1$ 与 $U_2$ 两个线丛，$V$ 与 $W$ 是 $E=V\oplus W$ 的两部分）可知第二式应为 $\bar\pi_*U_2=W$。译文照原文。]

\begin{prob}{10}
\textbf{抛物 Hitchin 纤维化。}如 [10] 所解释的，parahoric Higgs 丛的模空间，特别是上一章所定义的抛物 Higgs 丛的模空间，带有一个 Hitchin 映射；该映射是 Poisson 映射，其一般纤维是阿贝尔簇。借助文献 [10]，把抛物 Higgs 丛放在 parahoric Higgs 丛的框架下考虑：

\begin{enumerate}[label=(\roman*),leftmargin=2.2em]
\item 解释 parahoric Higgs 丛与抛物 Higgs 丛之间的关系；
\item 通过考虑 parahoric 整体幂零锥，用 $\operatorname{Bun}_{\cG}$ 的余切丛来描述抛物 $\cG$-Higgs 丛在 $0$ 上方的纤维。
\item 概述如何证明多重稳定抛物 Higgs 丛的模空间上的 Hitchin 映射是紧合的（proper）。
\item 考虑抛物 $\SL(4,\C)$-Higgs 丛，并比较 [72] 与 [10] 中所定义的 Hitchin 映射。
\end{enumerate}
\end{prob}

% [Tutor Note: p.25 问题 10 原文有三处小疵：标题 "Parabolic Hitchin fibration" 后脱句点（其余问题标题均有）；ii 中 "in terms the cotangent bundle" 脱介词（应为 in terms of）；iii 中 "summery" 应为 "summary"。译文按正确含义处理，正文不改原文句式。]
